





\section{What the Arc Reveals: Revisiting the Design Axes}
\label{sec:discussion:axes}

\cref{ch:framework} identified four design axes that govern dynamics-constrained motion generation: \textit{modeling burden}, \textit{dimensionality and factorization}, \textit{integration and pipeline brittleness}, and \textit{usability}.
Looking back across the thesis arc, each axis traces a coherent trajectory.

\paragraph{Modeling burden.}
\textit{PokeRRT} outsourced dynamics modeling entirely to a physics simulator, avoiding the need for analytical contact models at the cost of repeated simulation queries.
\textit{SFRRT*} replaced explicit liquid modeling with a learned classifier, trading model precision for coverage of complex geometries and obstacle configurations.
The payload diffusion model absorbed constraint satisfaction into the generation process itself: the ``model'' is the weights of the diffusion network, not a separate physics module.
Task-agnostic seeding takes this further, eliminating task-specific conditioning by representing constraints as samples from the constraint manifold directly.
The arc is a consistent progression from explicit simulation to specialized learned models to task-agnostic constraint representations.

\paragraph{Dimensionality and factorization.}
Kinodynamic planning in joint-velocity-acceleration space is fundamentally high-dimensional.
\textit{PokeRRT} manages this through random tree exploration; \textit{LazyBoE} manages it through pre-computation and lazy evaluation.
Diffusion-based methods sidestep the search problem entirely, compressing high-dimensional constraint spaces into learned latent representations and generating solutions directly.
The progression illustrates a shift from managing dimensionality through clever search to avoiding it through representation.

\paragraph{Integration and pipeline brittleness.}
Each kinodynamic system required carefully engineered components: a planner, a simulator, a failure model or liquid classifier, a time parameterizer, and their interfaces.
The diffusion-based systems collapsed multiple pipeline stages---the generator absorbs constraint checking, trajectory parameterization, and task adaptation simultaneously.
Task-agnostic seeding goes furthest: the same generator feeds into any downstream optimizer, eliminating the need for task-specific pipeline design entirely.

\paragraph{Usability.}
\textit{PokeRRT} required hand-tuning of planner parameters and scenario-specific design.
The whole-body and liquid systems required substantial per-task engineering.
Payload diffusion reduced the task specification to a single scalar.
\textit{DEFT} reduced it to an embodiment encoding.
Task-agnostic seeding reduces it to a set of constraint-satisfying state samples, which can be generated automatically from any robot dynamics model.
Each step along the axis reduces the expertise required to deploy a dynamics-constrained planner on a new task.